\section{Att kvadratkomplettera uttryck}

Metoden för att lösa andragradsekvationer som presenteras ovaan fungerar för att alla $x$-termer är samlade i en kvadrat, så som $(x - 13)^2$
Tack vare det kan vi ta alla okända termer ($x$-termer) på ena sidan i en ekvation och sedan ta kvadratroten ur båda led, vilket förvandlar andragradsekvationen till (oftast) två linjära ekvationer.

Andragradsuttryck där alla okända termer är samlade i en kvadrat kallas \emph{kvadratkompletterade}, och det är inte alla gånger som andragradsuttryck är skrivna på den formen.
I det här avsnittet ska vi lära oss att omvandla ett andragradsuttryck skrivet på allmän form ($ax^2 + bx + c$) till ett kvadratkompletterat uttryck.
Det arbetet kräver att vi är säkra på att lösa linjära ekvationer, men också att vi är bekväma med att multiplicera parentesuttryck som $(x+3)(2-x)$.

Hur hanterar vi sådana uttryck?

\subsection{Att multiplicera två parentesuttryck}

I matte 1 har vi tränat på att multiplicera parentesuttryck med ett tal.
Det gör vi genom att multiplicera \emph{varje term} i parentesen med talet framför parentesen:

\begin{center}
$7(2-x) = \framebox{7} \cdot 2 + \framebox{7} \cdot (-x) = 14 - 7x$

eller

$x(2y+x-3) = \framebox{x} \cdot 2y + \framebox{x} \cdot x + \framebox{x} \cdot (-3) = 2xy + x^2 - 3x$
\end{center}

Om vi inte har ett tal framför parentesen, utan istället ett helt parentesuttryck, kan vi göra på samma sätt:
Istället för att multiplicera varje term med $x$, multiplicerar vi med \emph{allt} som står framför parentesen.

Om vi exempelvis vill utveckla uttrycket $(x+3)(2-x)$ kan vi göra så här:

\smallskip
\begin{tabular}{l|p{5.7cm}}
  $\framebox{(x+3)} (2-x)$ & Vi ska multiplicera in $x+3$ i parentesen \\
  $= \framebox{(x+3)} \cdot 2 + \framebox{(x+3)} \cdot (-x)$ & Vi får nu två nya parentesuttryck, som vi måste utveckla vidare. Det går lättare att se vad vi gör om vi vänder på multiplikationen. \\
  $=2 \cdot \framebox{(x+3)} + (-x) \cdot \framebox{(x+3)}$ & Multiplicera in talen framför varje parentes, som vanligt. \\
  $=2 \cdot x + 2 \cdot 3 + (-x) \cdot x + (-x) \cdot 3$ & Snygga till uttrycken så att de går lättare att läsa. \\
  $=2x + 6 - x^2  - 3x$ & Förenkla uttrycket genom att slå samman $x$-termerna. (Men håll $x^2$ separat!) \\
  $=-x^2 - x + 6$ & \\
\end{tabular}
\smallskip

Den här metoden för att multiplicera parentesuttryck är bra för att den \emph{alltid} fungerar:
Om vi vill multiplicera längre uttryck eller krångligare uttryck kan vi alltid följa samma regel.
Men metoden är inte särskilt snabb, och därför kommer vi att lära oss ett par användbara genvägar.

\subsubsection{Genväg 1: Multiplicera alla möjliga termer}

Den vanligaste genvägen för att multiplicera parentesuttryck är att ta \emph{varje} term i den ena parentesen, multiplicera med \emph{varje} term i den andra, och summera resultaten.
När man lär sig denna metod brukar man ofta rita pilar mellan termerna, så att man ser att alla möjliga produkter bildats:

\begin{figure}
  \centering
  \includegraphics[width=0.7\textwidth]{bilder/parentesmultiplikation.svg}
  \caption{\label{fig:parentesmultiplikation}Utritning av pilar när man utvecklar uttrycket $(x+3)(x-2)$.}
\end{figure}

Exempel: Utveckla uttrycket $(x+3)(2-x)$.

\smallskip
\begin{tabular}{l|p{5.7cm}}
  $(x+3)(2-x)$ & Vi börjar med att multiplicera $x$ i första parentesen med varje term i andra parentesen. \\
  $= x \cdot 2 + x \cdot (-x) + (\cancel{x} + 3)(2-x)$ &  Vi fortsätter med att även multiplicera trean i första parentesen. \\
  $= x \cdot 2 + x \cdot (-x) + 3 \cdot 2 + 3 \cdot (-x)$ & Multiplicera in talen framför varje parentes, som vanligt. (Alla termer är nu multiplicerade.) \\
  $= 2x - x^2 + 6 - 3x$ & Förenkla uttrycket genom att slå samman $x$-termerna. \\
  $=-x^2 - x + 6$ & \\
\end{tabular}
\smallskip
